\IfFileExists{stacks-project.cls}{%
\documentclass{stacks-project}
}{%
\documentclass{amsart}
}

% For dealing with references we use the comment environment
\usepackage{verbatim}
\newenvironment{reference}{\comment}{\endcomment}
%\newenvironment{reference}{}{}
\newenvironment{slogan}{\comment}{\endcomment}
\newenvironment{history}{\comment}{\endcomment}

% For commutative diagrams we use Xy-pic
\usepackage[all]{xy}

% We use 2cell for 2-commutative diagrams.
\xyoption{2cell}
\UseAllTwocells

% We use multicol for the list of chapters between chapters
\usepackage{multicol}

% This is generally recommended for better output
\usepackage{lmodern}
\usepackage[T1]{fontenc}

% For cross-file-references
\usepackage{xr-hyper}

% Package for hypertext links:
\usepackage{hyperref}

% For any local file, say "hello.tex" you want to link to please
% use \externaldocument[hello-]{hello}
\externaldocument[introduction-]{introduction}
\externaldocument[conventions-]{conventions}
\externaldocument[sets-]{sets}
\externaldocument[categories-]{categories}
\externaldocument[topology-]{topology}
\externaldocument[sheaves-]{sheaves}
\externaldocument[sites-]{sites}
\externaldocument[stacks-]{stacks}
\externaldocument[fields-]{fields}
\externaldocument[algebra-]{algebra}
\externaldocument[brauer-]{brauer}
\externaldocument[homology-]{homology}
\externaldocument[derived-]{derived}
\externaldocument[simplicial-]{simplicial}
\externaldocument[more-algebra-]{more-algebra}
\externaldocument[smoothing-]{smoothing}
\externaldocument[modules-]{modules}
\externaldocument[sites-modules-]{sites-modules}
\externaldocument[injectives-]{injectives}
\externaldocument[cohomology-]{cohomology}
\externaldocument[sites-cohomology-]{sites-cohomology}
\externaldocument[dga-]{dga}
\externaldocument[dpa-]{dpa}
\externaldocument[sdga-]{sdga}
\externaldocument[hypercovering-]{hypercovering}
\externaldocument[schemes-]{schemes}
\externaldocument[constructions-]{constructions}
\externaldocument[properties-]{properties}
\externaldocument[morphisms-]{morphisms}
\externaldocument[coherent-]{coherent}
\externaldocument[divisors-]{divisors}
\externaldocument[limits-]{limits}
\externaldocument[varieties-]{varieties}
\externaldocument[topologies-]{topologies}
\externaldocument[descent-]{descent}
\externaldocument[perfect-]{perfect}
\externaldocument[more-morphisms-]{more-morphisms}
\externaldocument[flat-]{flat}
\externaldocument[groupoids-]{groupoids}
\externaldocument[more-groupoids-]{more-groupoids}
\externaldocument[etale-]{etale}
\externaldocument[chow-]{chow}
\externaldocument[intersection-]{intersection}
\externaldocument[pic-]{pic}
\externaldocument[weil-]{weil}
\externaldocument[adequate-]{adequate}
\externaldocument[dualizing-]{dualizing}
\externaldocument[duality-]{duality}
\externaldocument[discriminant-]{discriminant}
\externaldocument[derham-]{derham}
\externaldocument[local-cohomology-]{local-cohomology}
\externaldocument[algebraization-]{algebraization}
\externaldocument[curves-]{curves}
\externaldocument[resolve-]{resolve}
\externaldocument[models-]{models}
\externaldocument[functors-]{functors}
\externaldocument[equiv-]{equiv}
\externaldocument[pione-]{pione}
\externaldocument[etale-cohomology-]{etale-cohomology}
\externaldocument[proetale-]{proetale}
\externaldocument[relative-cycles-]{relative-cycles}
\externaldocument[more-etale-]{more-etale}
\externaldocument[trace-]{trace}
\externaldocument[crystalline-]{crystalline}
\externaldocument[spaces-]{spaces}
\externaldocument[spaces-properties-]{spaces-properties}
\externaldocument[spaces-morphisms-]{spaces-morphisms}
\externaldocument[decent-spaces-]{decent-spaces}
\externaldocument[spaces-cohomology-]{spaces-cohomology}
\externaldocument[spaces-limits-]{spaces-limits}
\externaldocument[spaces-divisors-]{spaces-divisors}
\externaldocument[spaces-over-fields-]{spaces-over-fields}
\externaldocument[spaces-topologies-]{spaces-topologies}
\externaldocument[spaces-descent-]{spaces-descent}
\externaldocument[spaces-perfect-]{spaces-perfect}
\externaldocument[spaces-more-morphisms-]{spaces-more-morphisms}
\externaldocument[spaces-flat-]{spaces-flat}
\externaldocument[spaces-groupoids-]{spaces-groupoids}
\externaldocument[spaces-more-groupoids-]{spaces-more-groupoids}
\externaldocument[bootstrap-]{bootstrap}
\externaldocument[spaces-pushouts-]{spaces-pushouts}
\externaldocument[spaces-chow-]{spaces-chow}
\externaldocument[groupoids-quotients-]{groupoids-quotients}
\externaldocument[spaces-more-cohomology-]{spaces-more-cohomology}
\externaldocument[spaces-simplicial-]{spaces-simplicial}
\externaldocument[spaces-duality-]{spaces-duality}
\externaldocument[formal-spaces-]{formal-spaces}
\externaldocument[restricted-]{restricted}
\externaldocument[spaces-resolve-]{spaces-resolve}
\externaldocument[formal-defos-]{formal-defos}
\externaldocument[defos-]{defos}
\externaldocument[cotangent-]{cotangent}
\externaldocument[examples-defos-]{examples-defos}
\externaldocument[algebraic-]{algebraic}
\externaldocument[examples-stacks-]{examples-stacks}
\externaldocument[stacks-sheaves-]{stacks-sheaves}
\externaldocument[criteria-]{criteria}
\externaldocument[artin-]{artin}
\externaldocument[quot-]{quot}
\externaldocument[stacks-properties-]{stacks-properties}
\externaldocument[stacks-morphisms-]{stacks-morphisms}
\externaldocument[stacks-limits-]{stacks-limits}
\externaldocument[stacks-cohomology-]{stacks-cohomology}
\externaldocument[stacks-perfect-]{stacks-perfect}
\externaldocument[stacks-introduction-]{stacks-introduction}
\externaldocument[stacks-more-morphisms-]{stacks-more-morphisms}
\externaldocument[stacks-geometry-]{stacks-geometry}
\externaldocument[moduli-]{moduli}
\externaldocument[moduli-curves-]{moduli-curves}
\externaldocument[examples-]{examples}
\externaldocument[exercises-]{exercises}
\externaldocument[guide-]{guide}
\externaldocument[desirables-]{desirables}
\externaldocument[coding-]{coding}
\externaldocument[obsolete-]{obsolete}
\externaldocument[fdl-]{fdl}
\externaldocument[index-]{index}

% Theorem environments.
%
\theoremstyle{plain}
\newtheorem{theorem}[subsection]{Theorem}
\newtheorem{proposition}[subsection]{Proposition}
\newtheorem{lemma}[subsection]{Lemma}

\theoremstyle{definition}
\newtheorem{definition}[subsection]{Definition}
\newtheorem{example}[subsection]{Example}
\newtheorem{exercise}[subsection]{Exercise}
\newtheorem{situation}[subsection]{Situation}

\theoremstyle{remark}
\newtheorem{remark}[subsection]{Remark}
\newtheorem{remarks}[subsection]{Remarks}

\numberwithin{equation}{subsection}

% Macros
%
\def\lim{\mathop{\mathrm{lim}}\nolimits}
\def\colim{\mathop{\mathrm{colim}}\nolimits}
\def\Spec{\mathop{\mathrm{Spec}}}
\def\Hom{\mathop{\mathrm{Hom}}\nolimits}
\def\Ext{\mathop{\mathrm{Ext}}\nolimits}
\def\SheafHom{\mathop{\mathcal{H}\!\mathit{om}}\nolimits}
\def\SheafExt{\mathop{\mathcal{E}\!\mathit{xt}}\nolimits}
\def\Sch{\mathit{Sch}}
\def\Mor{\mathop{\mathrm{Mor}}\nolimits}
\def\Ob{\mathop{\mathrm{Ob}}\nolimits}
\def\Sh{\mathop{\mathit{Sh}}\nolimits}
\def\NL{\mathop{N\!L}\nolimits}
\def\CH{\mathop{\mathrm{CH}}\nolimits}
\def\proetale{{pro\text{-}\acute{e}tale}}
\def\etale{{\acute{e}tale}}
\def\QCoh{\mathit{QCoh}}
\def\Ker{\mathop{\mathrm{Ker}}}
\def\Im{\mathop{\mathrm{Im}}}
\def\Coker{\mathop{\mathrm{Coker}}}
\def\Coim{\mathop{\mathrm{Coim}}}

% Boxtimes
%
\DeclareMathSymbol{\boxtimes}{\mathbin}{AMSa}{"02}

%
% Macros for moduli stacks/spaces
%
\def\QCohstack{\mathcal{QC}\!\mathit{oh}}
\def\Cohstack{\mathcal{C}\!\mathit{oh}}
\def\Spacesstack{\mathcal{S}\!\mathit{paces}}
\def\Quotfunctor{\mathrm{Quot}}
\def\Hilbfunctor{\mathrm{Hilb}}
\def\Curvesstack{\mathcal{C}\!\mathit{urves}}
\def\Polarizedstack{\mathcal{P}\!\mathit{olarized}}
\def\Complexesstack{\mathcal{C}\!\mathit{omplexes}}
% \Pic is the operator that assigns to X its picard group, usage \Pic(X)
% \Picardstack_{X/B} denotes the Picard stack of X over B
% \Picardfunctor_{X/B} denotes the Picard functor of X over B
\def\Pic{\mathop{\mathrm{Pic}}\nolimits}
\def\Picardstack{\mathcal{P}\!\mathit{ic}}
\def\Picardfunctor{\mathrm{Pic}}
\def\Deformationcategory{\mathcal{D}\!\mathit{ef}}
\usepackage{cleveref}
\crefname{originallemma}{original lemma}{original lemmas}
\Crefname{originallemma}{Original lemma}{Original lemmas}
\makeatletter
\newlabel{lemma-relative-mayer-vietoris}{{0GKN}{source}{}{sourcepointer-1}{}}
\newlabel{lemma-relative-mayer-vietoris@cref}{{[originallemma][0][]0GKN}{[1][0][]source}{}{}{}}
\makeatother
\crefname{originallemma}{original lemma}{original lemmas}
\Crefname{originallemma}{Original lemma}{Original lemmas}
\makeatletter
\newlabel{lemma-shriek-composition}{{0F7J}{source}{}{sourcepointer-2}{}}
\newlabel{lemma-shriek-composition@cref}{{[originallemma][0][]0F7J}{[1][0][]source}{}{}{}}
\makeatother
\crefname{originallemma}{original lemma}{original lemmas}
\Crefname{originallemma}{Original lemma}{Original lemmas}
\makeatletter
\newlabel{lemma-qf-f-shriek-constructible}{{0GK3}{source}{}{sourcepointer-3}{}}
\newlabel{lemma-qf-f-shriek-constructible@cref}{{[originallemma][0][]0GK3}{[1][0][]source}{}{}{}}
\makeatother
\crefname{originallemma}{original lemma}{original lemmas}
\Crefname{originallemma}{Original lemma}{Original lemmas}
\makeatletter
\newlabel{lemma-derived-lower-shriek-bounded}{{0G2A}{source}{}{sourcepointer-4}{}}
\newlabel{lemma-derived-lower-shriek-bounded@cref}{{[originallemma][0][]0G2A}{[1][0][]source}{}{}{}}
\makeatother
\crefname{originallemma}{original lemma}{original lemmas}
\Crefname{originallemma}{Original lemma}{Original lemmas}
\makeatletter
\newlabel{derived-lemma-two-ss-complex-functor}{{015J}{source}{}{sourcepointer-5}{}}
\newlabel{derived-lemma-two-ss-complex-functor@cref}{{[originallemma][0][]015J}{[1][0][]source}{}{}{}}
\makeatother
\crefname{originallemma}{original lemma}{original lemmas}
\Crefname{originallemma}{Original lemma}{Original lemmas}
\makeatletter
\newlabel{lemma-constructible-shriek-rel-dim-1}{{0GKZ}{source}{}{sourcepointer-6}{}}
\newlabel{lemma-constructible-shriek-rel-dim-1@cref}{{[originallemma][0][]0GKZ}{[1][0][]source}{}{}{}}
\makeatother
\title{Constructibility of the cohomology sheaves of Rf\_! for separated finitely presented morphisms}
\author{The Stacks Project Authors}
\date{stacks-project:more-etale.tex:local-snapshot-20261012; GNU FDL 1.2 or later}
\begin{document}
\maketitle
\setcounter{subsection}{1}
\noindent\textit{Editorial note.} All delivered mathematical text is the author's own, copied from the pinned Stacks Project chapter file more-etale.tex; no statement or proof step was rewritten, reordered or removed. The counted span is the author environment and its own proof environment. References inside the span are handed to the assembler as navigation stubs that carry the canonical tag of the cited tag only; whether each of them is load-bearing is recorded on the stub itself.\par\medskip
\noindent\textit{Editorial note.} Editorial note: the final reduction uses the first spectral sequence of the two-spectral-sequence lemma for a complex functor as a named input (the author also offers an alternative truncation argument); the per-term bookkeeping of the spectral sequence reduction is not written out.\par\medskip
\noindent\textit{Allowed standard prerequisites (editorial).}\begin{itemize}
\item Relative Mayer-Vietoris for lower shriek (Lemma lemma-relative-mayer-vietoris, canonical tag 0GKN). cited at proof line 4940. Together with the induction principle it reduces the claim from general X to affine X.
\item Composition of lower shriek (Lemma lemma-shriek-composition, canonical tag 0F7J). cited at proof lines 4947-4948. Gives Rf\_! = Rp\_! after Ri\_! for the factorization through the closed immersion.
\item Constructibility for locally quasi-finite finitely presented morphisms (Lemma lemma-qf-f-shriek-constructible, canonical tag 0GK3). cited at proof lines 4949-4950. Handles the closed immersion part since Ri\_! = i\_! for a locally quasi-finite immersion.
\item Boundedness of the derived lower shriek functors (Lemma lemma-derived-lower-shriek-bounded, canonical tag 0G2A). cited at proof lines 4961-4964. Rf\_! has finite cohomological dimension, allowing K to be assumed bounded below.
\item Two spectral sequences for a complex functor (Derived Categories, Lemma derived-lemma-two-ss-complex-functor, canonical tag 015J). cited at proof lines 4964-4966. Its first spectral sequence reduces the claim from bounded below K to cohomology sheaves of K.
\item Constructibility for relative dimension 1 (Lemma lemma-constructible-shriek-rel-dim-1, canonical tag 0GKZ). cited at proof lines 4967-4968. The A\textasciicircum{}1 case target at which the reduction arrives.
\end{itemize}\medskip
\ifdefined\hypertarget\hypertarget{sourcepointer-1}{}\fi
\noindent\textit{Source reference.} Original lemma 0GKN: sources/stacks-chow/tags.txt, source lines 20693--20693. Cited at proof line 4940: together with the induction principle it reduces the claim to affine X.\par\medskip
\ifdefined\hypertarget\hypertarget{sourcepointer-2}{}\fi
\noindent\textit{Source reference.} Original lemma 0F7J: sources/stacks-chow/tags.txt, source lines 19009--19009. Cited at proof lines 4947-4948: gives Rf\_! = Rp\_! after Ri\_! for the factorization through the closed immersion.\par\medskip
\ifdefined\hypertarget\hypertarget{sourcepointer-3}{}\fi
\noindent\textit{Source reference.} Original lemma 0GK3: sources/stacks-chow/tags.txt, source lines 20673--20673. Cited at proof lines 4949-4950: handles the closed immersion part since Ri\_! = i\_! for a locally quasi-finite immersion.\par\medskip
\ifdefined\hypertarget\hypertarget{sourcepointer-4}{}\fi
\noindent\textit{Source reference.} Original lemma 0G2A: sources/stacks-chow/tags.txt, source lines 20050--20050. Cited at proof lines 4961-4964: Rf\_! has finite cohomological dimension, so K may be assumed bounded below.\par\medskip
\ifdefined\hypertarget\hypertarget{sourcepointer-5}{}\fi
\noindent\textit{Source reference.} Original lemma 015J: sources/stacks-chow/tags.txt, source lines 1510--1510. Cited at proof lines 4964-4966: its first spectral sequence reduces the claim from bounded below K to the cohomology sheaves of K.\par\medskip
\ifdefined\hypertarget\hypertarget{sourcepointer-6}{}\fi
\noindent\textit{Source reference.} Original lemma 0GKZ: sources/stacks-chow/tags.txt, source lines 20704--20704. Cited at proof lines 4967-4968: the A\textasciicircum{}1 case target at which the reduction arrives.\par\medskip
\begin{theorem}
\label{theorem-constructible-shriek}
Let $f : X \to Y$ be a separated morphism of finite presentation
of quasi-compact and quasi-separated schemes. Let $\Lambda$ be a
Noetherian ring. Let $K$ be an object of $D^+_{tors, c}(X_\etale, \Lambda)$
or of $D_c(X_\etale, \Lambda)$ in case $\Lambda$ is torsion.
Then $Rf_!K$ has constructible cohomology sheaves, i.e.,
$Rf_!K$ is in $D^+_{tors, c}(Y_\etale, \Lambda)$
or in $D_c(Y_\etale, \Lambda)$ in case $\Lambda$ is torsion.
\end{theorem}

\begin{proof}
The question is local on $Y$ hence we may and do assume $Y$ is affine.
By the induction principle and Lemma \ref{lemma-relative-mayer-vietoris}
we reduce to the case where $X$ is also affine.

\medskip\noindent
Assume $X$ and $Y$ are affine. Since $X$ is of finite presentation,
we can choose a closed immersion $i : X \to \mathbf{A}^n_Y$ which is
of finite presentation. If $p : \mathbf{A}^n_Y \to Y$ denotes
the structure morphism, then we see that $Rf_! = Rp_! \circ Ri_!$
by Lemma \ref{lemma-shriek-composition}. By
Lemma \ref{lemma-qf-f-shriek-constructible}
we have the result for $Ri_! = i_!$. Hence we may assume $f$
is the projection morphism $\mathbf{A}^n_Y \to Y$.
Since we can view $f$ as the composition
$$
X = \mathbf{A}^n_Y \to \mathbf{A}^{n - 1}_Y \to
\mathbf{A}^{n - 2}_S \to \ldots \to
\mathbf{A}^1_Y \to Y
$$
we may assume $n = 1$.

\medskip\noindent
Assume $Y$ is affine and $X = \mathbf{A}^1_Y$. Since $Rf_!$ has
finite cohomological dimension
(Lemma \ref{lemma-derived-lower-shriek-bounded})
we may assume $K$ is bounded below. Using the first spectral sequence of
Derived Categories, Lemma \ref{derived-lemma-two-ss-complex-functor}
(or alternatively using an argument with truncations),
we reduce to showing the result of
Lemma \ref{lemma-constructible-shriek-rel-dim-1}.
\end{proof}

\end{document}
